\section{Discussion and Limitations}
\label{sec:discussion}

\subsection{Validity of the Results}
Four limitations apply to the results of \cref{sec:results} and determine how far they may
be generalised.

\paragraph{Simulation only.} Every measurement in this report was taken in the simulation.
None was taken on the printed robot, for which there is no recorded data beyond the
photographs and the outcome of the two competition rounds. The simulated motors are ideal
within their limits, the floor is flat and evenly lit, and the infrared model has no
dependence on ambient light. The results show that the software is correct with respect to
its own model; they do not predict the performance of the hardware.

\paragraph{A model built on CAD values.} The mass properties on which \cref{sec:mass} rests
are those that the CAD program calculated. The robot was not weighed, and the motor data in
the model are placeholders. The conclusions drawn there are large enough to survive
moderate errors, since a wheel load of 18\,\% would have to be wrong by a factor of several
before the conclusion changed, but the figures are not measurements.

\paragraph{A favourable arena.} The arena was designed together with the software. Its
markings have the widths that the array needs, its walls have the textures that the SLAM
system needs, and its lighting casts no shadows. A competition field offers none of these
by arrangement. The arena shows that the approach works when its assumptions hold, and
\cref{sec:arena} states what those assumptions are.

\paragraph{A small number of runs.} The mission was run twelve times and each line-following
condition once. The runs differ in the order of the colours and in the obstacles, not in
the start pose, the map or the lighting, and the map was the same in all of them.
Differences of a few millimetres or a few seconds between conditions are within what a
repetition could change.

\subsection{Known Issues}
\label{sec:issues}
\begin{description}
  \item[The wheels are lightly loaded.] The model puts 82\,\% of the weight on the caster
        (\cref{sec:mass}). The simulation works within the resulting traction limit only
        because the drive is limited to 1.0\unit{m/s^2}. On the hardware, the caster should
        be moved towards the gripper and the camera to the centre line, after the robot has
        been weighed on its three supports.
  \item[The caster does not reach the floor in CAD.] The dome ends 0.74\unit{mm} above the
        wheel contact plane (\cref{sec:urdf}). The simulation corrects this in its collision
        geometry; the physical part needs the same correction.
  \item[The jaws do not meet on the centre line.] The second jaw is $7.25^\circ$ off the
        mirror image of the first in the exported assembly (\cref{sec:gripper}).
  \item[The footprint used for planning is a disc.] The navigation treats the robot as a
        disc of 0.13\unit{m} radius about the midpoint of the axle. The closed jaws reach
        0.21\unit{m} ahead of that point and lie outside the disc. The inflation of the
        costmap by 0.25\unit{m} covers the difference in the present arena; a polygonal
        footprint would be the correct description.
  \item[Low obstacles are not seen.] Anything lower than 3\unit{cm} is treated as floor, and
        nothing closer than 0.15\unit{m} in front of the robot is seen at all.
  \item[Localisation depends on a threshold.] The condition that rejects look-alike views
        (\cref{sec:slam}) is a single number chosen for this arena and this camera height.
        It rejects some true matches, and a different arena would need it chosen again.
  \item[Dimensions are repeated in three files.] The station poses and the drop squares
        mirror the arena file by hand (\cref{sec:arena}).
  \item[No feedback from the jaws.] The mission assumes that a coin is caged when the array
        is on the crossbar and the jaws have had time to close. It cannot tell that a coin
        was missed or lost on the way. Only the referee sees this.
  \item[One line sensor, no fallback.] This is the failure of the state round
        (\cref{sec:electronics}), and the simulated mission has the same dependency: every
        approach ends on the infrared array.
\end{description}

\subsection{What I Would Do Differently}
Three changes follow from this work, in order of value.
\begin{enumerate}
  \item \emph{Check the mass properties during the design, not after it.} The calculation of
        \cref{sec:mass} takes minutes once the CAD model exists. Done during the fifteen
        days, it would have moved the caster and the camera before the first print.
  \item \emph{Treat the line sensor as a part that can fail.} A check of all eight channels
        on a black and a white reference at power-up, a soldered harness, and a degraded
        mode that crosses the line section on odometry and vision are the measures
        that address a failure of this kind.
  \item \emph{Build the simulation first.} The simulation was written after the competition.
        Written before it, it would have allowed the mission software to be developed
        while the parts were printing.
\end{enumerate}

\subsection{Future Work}
\begin{itemize}
  \item Weigh the printed robot on its supports, measure the motors, and replace the
        placeholder values of the model.
  \item Move the caster and the camera as proposed and repeat the analysis of \cref{sec:mass}.
  \item Run the mission on the hardware, beginning with the line follower and the docking,
        which need only the infrared array and the encoders.
  \item Add a contact or current signal from the gripper so that a missed coin is detected.
  \item Derive the station poses and the referee squares from the arena file.
  \item Vary the start pose, the lighting and the floor friction in the simulation, and
        repeat each condition several times.
\end{itemize}
